% Shared Somerville styling and MetaFun graphics for the book transcriptions.

\startenvironment mechanism-environment

\usepath[..]
\input type-imp-somerville.mkxl

\setupbodyfontenvironment[default][em=italic]
\setupSomerville
  {Width=100,Weight=400,Optical Size=12,Math Weight=700,Arrow Length=200,Arrow Head=100,Integral Slant=100}
  {somerville-heavens}
\setupbodyfont[somerville-heavens,11pt]

% Base size matches cap heights: Playfair 708/1240 at 11pt, Porson 1440/2048.
% Base Porson size = 11 * (708/1240) / (1440/2048) = 8.932473pt.
% Enlarge Greek by 10% from that base, including the varpi fallback below.
% Script sizes use 75% and 55% of the text size.
% Use BMP Greek codes: Porson does not encode mathematical-alphanumeric Greek.
% ConTeXt's weight effect adds 0.008em to stroke thickness: about 10% of
% Porson's typical 160/2048em stem. It scales with each font size.
\definefontfeature[heavens-greek-weight][default]
  [effect={effect=both,weight=0.08}]
\definefont[HeavensSymbols]            [file:GFSPorson.ttf*heavens-greek-weight at 9.82575pt]
\definefont[HeavensScriptSymbols]      [file:GFSPorson.ttf*heavens-greek-weight at 7.3693125pt]
\definefont[HeavensScriptScriptSymbols][file:GFSPorson.ttf*heavens-greek-weight at 5.4041625pt]
\unexpanded\def\HeavensGreek#1{%
  \mathord{\mathchoice
    {\hbox{\HeavensSymbols #1}}%
    {\hbox{\HeavensSymbols #1}}%
    {\hbox{\HeavensScriptSymbols #1}}%
    {\hbox{\HeavensScriptScriptSymbols #1}}}%
}
\unexpanded\def\HeavensSigma  {\HeavensGreek{\char"03A3\relax}}
\unexpanded\def\HeavensAlpha  {\HeavensGreek{\char"03B1\relax}}
\unexpanded\def\HeavensBeta   {\HeavensGreek{\char"03B2\relax}}
\unexpanded\def\HeavensGamma  {\HeavensGreek{\char"03B3\relax}}
\unexpanded\def\HeavensDelta  {\HeavensGreek{\char"03B4\relax}}
\unexpanded\def\HeavensEpsilon{\HeavensGreek{\char"03B5\relax}}
\unexpanded\def\HeavensTheta  {\HeavensGreek{\char"03B8\relax}}
\unexpanded\def\HeavensLambda {\HeavensGreek{\char"03BB\relax}}
\unexpanded\def\HeavensMu     {\HeavensGreek{\char"03BC\relax}}
\unexpanded\def\HeavensNu     {\HeavensGreek{\char"03BD\relax}}
\unexpanded\def\HeavensPi     {\HeavensGreek{\char"03C0\relax}}
\unexpanded\def\HeavensRho    {\HeavensGreek{\char"03C1\relax}}
\unexpanded\def\HeavensPhi    {\HeavensGreek{\char"03C6\relax}}
\unexpanded\def\HeavensPsi    {\HeavensGreek{\char"03C8\relax}}
\unexpanded\def\HeavensOmega  {\HeavensGreek{\char"03C9\relax}}

% Porson lacks the mathematical pi symbol (varpi); retain its distinct shape.
\definefont[HeavensVariantSymbols]            [file:latinmodern-math.otf*heavens-greek-weight at 12.1pt]
\definefont[HeavensScriptVariantSymbols]      [file:latinmodern-math.otf*heavens-greek-weight at 9.075pt]
\definefont[HeavensScriptScriptVariantSymbols][file:latinmodern-math.otf*heavens-greek-weight at 6.655pt]
\unexpanded\def\HeavensVarpi{%
  \mathord{\mathchoice
    {\hbox{\HeavensVariantSymbols\char"1D71B\relax}}%
    {\hbox{\HeavensVariantSymbols\char"1D71B\relax}}%
    {\hbox{\HeavensScriptVariantSymbols\char"1D71B\relax}}%
    {\hbox{\HeavensScriptScriptVariantSymbols\char"1D71B\relax}}}%
}
\unexpanded\def\HeavensCircle{\circ}

\definepapersize[heavens][width=150mm,height=240mm]
\setuppapersize[heavens][heavens]
% Aged paper: change strength to 0 for flat cream, or increase it for more age.
% The seed and page number make the texture repeatable and different per page.
\definecolor[heavens-paper][h=FBF4D8]
\def\HeavensPaperStrength{0}
\def\HeavensPaperSeed{1831}

\startluacode
do
    local exp, sin, floor, min, max = math.exp, math.sin, math.floor, math.min, math.max
    local function byte(value)
        return max(0, min(255, value))
    end

    function MP.HeavensPaperColor(value, x, y)
        -- These dimensions match the MetaFun bytemap below (about 122 dpi).
        local u, v = x / 719, y / 1151
        local page = MP.heavenspaperpage or 1
        local seed = MP.heavenspaperseed or 1831
        local strength = MP.heavenspaperstrength or 1

        -- Large cloudy changes in the paper's yellow/brown tone.
        local cloud = (value - 127.5) / 127.5

        -- Soft edge ageing, with a slightly wider shadow at the binding.
        -- The chapter begins on printed page 4, with binding on the right.
        local edge = .65 * exp(-min(u, 1-u) / .018)
                   + .45 * exp(-min(v, 1-v) / .015)
        local bindingdistance = page % 2 == 1 and (1-u) or u
        local binding = exp(-bindingdistance / .045)

        -- Deterministic, very low contrast grain; no global random state.
        local hash = sin(x * 12.9898 + y * 78.233 + page * 37.719 + seed * .123) * 43758.5453
        local grain = 2 * (hash - floor(hash)) - 1

        return byte(251 + strength * (-8  * cloud - 3  * edge - 3.5 * binding + .5 * grain)),
               byte(244 + strength * (-12 * cloud - 6  * edge - 7   * binding + .8 * grain)),
               byte(216 + strength * (-24 * cloud - 10 * edge - 11  * binding +      grain))
    end
end
\stopluacode

\startuseMPgraphic{heavens:aged-paper}
  draw lmt_noise [
    bytemap       = 31,
    nx            = 720,
    ny            = 1152,
    nz            = 3,
    method        = "perlin",
    iterations    = 3,
    frequency     = .008,
    persistence   = .45,
    lacunarity    = 2,
    minimum       = 0,
    maximum       = 255,
    z             = \HeavensPaperSeed/100 + .73*\the\realpageno,
    colorfunction = "HeavensPaperColor",
  ] xysized (OverlayWidth,OverlayHeight);
\stopuseMPgraphic

\defineoverlay[heavens:aged-paper]
  [{\ctxlua{
      MP.heavenspaperpage = \the\realpageno;
      MP.heavenspaperseed = \HeavensPaperSeed;
      MP.heavenspaperstrength = \HeavensPaperStrength;
    }\useMPgraphic{heavens:aged-paper}}]
\setupbackgrounds[page]
  [background=heavens:aged-paper,backgroundcolor=heavens-paper]
\setuplayout
  [backspace=17mm,topspace=10mm,width=116mm,height=218mm,
   header=7mm,headerdistance=4mm,footer=0mm]
\setuppagenumbering[location=]
\setupinterlinespace[line=12.8pt]
\setupwhitespace[none]
\setupindenting[no]
\setupalign[normal,verytolerant]
\setupformulas[spacebefore=3pt,spaceafter=3pt]
\setupfloats[spacebefore=none,spaceafter=none]

% Upright force and point labels, as in the original; variables stay italic.
\unexpanded\def\Force{\mathrm{F}}
\unexpanded\def\ForceX{\mathrm{X}}
\unexpanded\def\ForceY{\mathrm{Y}}
\unexpanded\def\ForceZ{\mathrm{Z}}
\unexpanded\def\Resultant{\Force_{1}}
\unexpanded\def\Der#1#2{\frac{\partial #1}{\partial #2}}
\unexpanded\def\SumForce{\HeavensSigma\mathpunct{.}\Force\mathpunct{.}}
\unexpanded\def\Article#1{\par\noindent\hskip1em\relax #1.\enspace}
\unexpanded\def\HeavensChapterHeaders#1#2{%
  \setupheadertexts
    [{\hbox to \makeupwidth{%
       \ifodd\userpageno
         \rlap{Chap. #1.]}\hss{#2}\hss\llap{\userfolio}%
       \else
         \rlap{\userfolio}\hss{#2}\hss\llap{[Book I.}%
       \fi}}]%
}

\startuseMPgraphic{heavens:figure10}
  numeric u; u := .20mm;
  pair O,A,B,C,P,Q,R,M;
  O := (0,0); A := (68u,0); B := (-18u,-22u); C := (0,52u);
  P := A+B; Q := A+C; R := B+C; M := A+B+C;
  pickup pencircle scaled .35pt;
  draw O--A--P--B--cycle;
  draw C--Q--M--R--cycle;
  draw O--C; draw A--Q; draw B--R; draw P--M;
  draw O--M; draw A--M; draw B--M; draw C--M;
  draw O--(116u,0); draw O--(0,92u); draw O--2.4B;
  label.rt(textext("\tfx\m{x}"),(116u,0));
  label.top(textext("\tfx\m{z}"),(0,92u));
  label.bot(textext("\tfx\m{y}"),2.4B);
  label.llft(textext("\tfx\m{o}"),O);
  label.rt(textext("\tfx A"),A);
  label.bot(textext("\tfx B"),B);
  label.ulft(textext("\tfx C"),C);
  label.bot(textext("\tfx P"),P);
  label.urt(textext("\tfx Q"),Q);
  label.lft(textext("\tfx R"),R);
  label.ulft(textext("\tfx\m{m}"),M);
  label(textext("\tfx{\it fig.}\,10."),(68u,98u));
\stopuseMPgraphic

\startuseMPgraphic{heavens:figure11}
  numeric u; u := .20mm;
  pair E,O,A,B,C,P,Q,R,M;
  E := (-46u,-46u);
  O := (0,0); A := (65u,0); B := (-23u,-17u); C := (0,44u);
  P := A+B; Q := A+C; R := B+C; M := A+B+C;
  pickup pencircle scaled .35pt;
  draw E--(112u,ypart E);
  draw E--(xpart E,91u);
  draw E--(E+(-53u,-38u));
  draw O--A--P--B--cycle;
  draw C--Q--M--R--cycle;
  draw O--C; draw A--Q; draw B--R; draw P--M;
  draw O--M dashed evenly scaled .5;
  label.rt(textext("\tfx\m{x}"),(112u,ypart E));
  label.top(textext("\tfx\m{z}"),(xpart E,91u));
  label.bot(textext("\tfx\m{y}"),E+(-53u,-38u));
  label.ulft(textext("\tfx\m{o'}"),E);
  label.llft(textext("\tfx\m{o}"),O);
  label.rt(textext("\tfx A"),A);
  label.bot(textext("\tfx B"),B);
  label.ulft(textext("\tfx C"),C);
  label.bot(textext("\tfx P"),P);
  label.urt(textext("\tfx Q"),Q);
  label.lft(textext("\tfx R"),R);
  label.lrt(textext("\tfx\m{m}"),M);
  label(textext("\tfx{\it fig.}\,11."),(80u,78u));
\stopuseMPgraphic


% Fit a long display without allowing a break inside an equation.
\unexpanded\def\FitFormula#1{%
  \startformula
    \hbox{\scale[maxwidth=\textwidth]{\m{\displaystyle #1}}}%
  \stopformula
}
\stopenvironment
